\section{Proofs and Auxiliary Results}
\label{app:theory}
\label{app:learning_proof}
\setcounter{equation}{0}

\subsection{Maximum-Sharpe direction}
\label{app:maximum_sharpe}

\begin{proof}[Proof of Proposition~\ref{prop:population_optimal_policy}]
\textit{Step 1: separate direction and scale.}
Every nonzero $W\in\Hh$ can be written as $W=aU$, with
$\|U\|_{\Hh}=1$ and $a\in\mathbb R$. Conversely, every such $aU$ belongs
to $\Hh$, and $a=0$ includes the zero policy. Thus this parametrization does
not change the feasible set. Since the payoff is linear in the policy,
\[
\min_{W\in\Hh}Q(W)
=
\min_{\substack{U\in\Hh\\\|U\|_{\Hh}=1}}
\min_{a\in\mathbb R}
\E\!\left[\left(1-aR^p_{t+1}(U)\right)^2\right].
\]
A zero-payoff direction has loss one for every scale. A nonzero constant
payoff would give minimum loss zero after rescaling. The assumption
$0<Q(W^\star)<1$ therefore allows zero-variance directions to be omitted.

\medskip
\noindent\textit{Step 2: optimize the scale.}
Fix a direction $U$ with positive payoff variance and write
\[
m_U:=\E[R^p_{t+1}(U)],
\qquad
s_U^2:=\E[(R^p_{t+1}(U))^2]>0.
\]
Expanding the square and differentiating gives
\[
Q(aU)=1-2am_U+a^2s_U^2,
\qquad
\frac{\partial Q(aU)}{\partial a}=-2m_U+2as_U^2.
\]
The second derivative is $2s_U^2>0$, so the unique minimizing scale is
$a^\star(U)=m_U/s_U^2$. Substitution yields
\[
\begin{aligned}
\min_{a\in\mathbb R}Q(aU)
&=1-2\frac{m_U}{s_U^2}m_U
 +\left(\frac{m_U}{s_U^2}\right)^2s_U^2\\
&=1-\frac{m_U^2}{s_U^2}\\
&=\frac{\Var(R^p_{t+1}(U))}
        {\Var(R^p_{t+1}(U))+m_U^2}\\
&=\frac{1}{1+\SR^2(R^p_{t+1}(U))}.
\end{aligned}
\]

\medskip
\noindent\textit{Step 3: compare directions.}
Consequently,
\[
\min_{W\in\Hh}Q(W)
=
\min_{\substack{U\in\Hh\\\|U\|_{\Hh}=1}}
\frac{1}{1+\SR^2(R^p_{t+1}(U))}.
\]
The right-hand side is minimized by a maximum-squared-Sharpe direction,
because $x\mapsto(1+x)^{-1}$ is strictly decreasing on $[0,\infty)$.
The corresponding optimally scaled payoff has mean
$m_U^2/s_U^2>0$, since the minimum is strictly below one. It therefore
attains the positive maximum Sharpe ratio, not its negative orientation.
\end{proof}

\subsection{Pricing-kernel interpretation and Hansen--Jagannathan equality}
\label{app:sdf_hj}

\begin{proof}[Proof of Proposition~\ref{prop:sdf_hj}]
Fix an economy and suppress its subscript in this proof.

\medskip
\noindent\textit{Step 1: portfolio orthogonality.}
For any $V\in\Hh$, expansion of the quadratic criterion gives
\[
\begin{aligned}
Q(W^\star+aV)
&=Q(W^\star)
 -2a\E\!\left[
  (1-R^p_{t+1}(W^\star))R^p_{t+1}(V)
 \right]\\
&\quad+a^2\E[(R^p_{t+1}(V))^2].
\end{aligned}
\]
Optimality at $a=0$ implies
$\E[m^\star_{t+1}R^p_{t+1}(V)]=0$ for every $V$. In particular,
choosing $V=W^\star$ gives
\[
\E[R^p_{t+1}(W^\star)]
=
\E[(R^p_{t+1}(W^\star))^2].
\]
Since $m^\star=1-R^p(W^\star)$, it follows that
\[
\E[m^\star]
=\E[(m^\star)^2]
=Q(W^\star)
=:q^\star\in(0,1).
\]
Proposition~\ref{prop:population_optimal_policy} then gives
$q^\star=1/(1+(\SR^\star)^2)$.

\medskip
\noindent\textit{Step 2: normalize the pricing kernel.}
For $\widetilde m=m^\star/q^\star$,
\[
\E[\widetilde m]=1,
\qquad
\E[\widetilde m^2]=\frac{1}{q^\star},
\qquad
\Var(\widetilde m)=\frac{1-q^\star}{q^\star}=(\SR^\star)^2.
\]

\medskip
\noindent\textit{Step 3: verify the minimum-variance property.}
Let $d$ be any square-integrable, possibly signed, random variable with
$\E[d]=1$ and $\E[dR^p(V)]=0$ for every $V\in\Hh$.
Then
\[
\E[dm^\star]
=\E[d]-\E[dR^p(W^\star)]
=1.
\]
Hence $\E[d\widetilde m]=1/q^\star=\E[\widetilde m^2]$, and therefore
\[
\Var(d)
=
\Var(\widetilde m)+\E[(d-\widetilde m)^2]
\geq (\SR^\star)^2.
\]
Thus $\widetilde m$ attains the Hansen--Jagannathan bound on the specified
managed excess-return space, with no positivity restriction. None of these
steps asserts the stronger conditional restriction in
Condition~\ref{ec:pricing}.
\end{proof}

\subsection{Managed-portfolio Markowitz representation}
\label{app:managed_markowitz}

\begin{proof}[Proof of Proposition~\ref{prop:managed_markowitz}]
Under the integrability conditions verified in
Lemma~\ref{lem:primitive_consequences}, the reproducing property gives
$R^p_{t+1}(W)=\langle W,X_t\rangle_{\Hh_K}$. Therefore
\[
Q(W)
=1-2\langle W,\mu_{\Hh_K}\rangle_{\Hh_K}
 +\langle W,\Sigma_{\Hh_K}W\rangle_{\Hh_K}.
\]
For an arbitrary direction $V\in\Hh_K$,
\[
DQ(W)[V]
=2\langle\Sigma_{\Hh_K}W-\mu_{\Hh_K},V\rangle_{\Hh_K}.
\]
The population first-order condition is consequently
\[
\Sigma_{\Hh_K}W_K^\star=\mu_{\Hh_K}.
\]
If the inverse is defined on $\mu_{\Hh_K}$ in the identifiable subspace,
this gives the population representation in the proposition. The normal
equation itself does not require a bounded inverse on all of $\Hh_K$.

The empirical penalized criterion is
\[
1-2\langle W,\widehat\mu_T\rangle_{\Hh_K}
 +\langle W,\widehat\Sigma_TW\rangle_{\Hh_K}
 +\lambda\|W\|_{\Hh_K}^2.
\]
Its derivative vanishes exactly when
\[
(\widehat\Sigma_T+\lambda I)\widehat W_{\lambda,T}=\widehat\mu_T.
\]
Since $\widehat\Sigma_T$ is positive and $\lambda>0$, the operator on the
left is boundedly invertible, with inverse norm at most $1/\lambda$.
This yields the unique empirical policy in
\eqref{eq:markowitz_representation}.
\end{proof}

\subsection{Consequences of the primitive conditions}
\label{app:primitive_consequences}

\begin{lemma}[Consequences of return-tail regularity]
\label{lem:primitive_consequences}
Fix $\varepsilon\in\mathcal E$, and put
$C_{\rm sc}=\sqrt{8K_4(1+K_4)}$. Write $\Sigma=\Sigma_{\Hh_K,\varepsilon}$.
\begin{enumerate}[label=\textup{(\roman*)},leftmargin=*,itemsep=0.35em]
\item The feature $X_t=N^{-1}\sum_iR_{i,t+1}K(\cdot,Z_{i,t})$ is strongly
measurable,
\[
R^p_{t+1}(W)=\langle W,X_t\rangle_{\Hh_K},\qquad
\E_\varepsilon\!\left[
\left.
\exp\!\left(\frac{\|X_t\|_{\Hh_K}^2}{B_X^2}\right)
\right|\mathscr F_t
\right]\leq2
\quad\text{a.s.}
\]
In particular,
\[
\mathbb P_\varepsilon(\|X_t\|>u\mid\mathscr F_t)
\leq2e^{-u^2/B_X^2}
\quad\text{a.s.},
\]
and $\E_\varepsilon\|X_t\|^{2q}\leq2q!B_X^{2q}$ for integers $q\geq1$.
These are moment bounds, not almost-sure bounds.

\item The operator $\Sigma$ is positive, self-adjoint, trace class, and compact,
with
\[
\Tr\Sigma=\E_\varepsilon\|X_t\|^2\leq B_X^2\log2.
\]

\item Writing $\mu=\E_\varepsilon X_t$, the minimum-norm target satisfies
$\Sigma W^\star_\varepsilon=\mu$. For every $W\in\Hh_K$,
\[
Q_\varepsilon(W)-Q_\varepsilon(W^\star_\varepsilon)
=\|\Sigma^{1/2}(W-W^\star_\varepsilon)\|_{\Hh_K}^2.
\]

\item Let $P_t^\star=R^p_{t+1}(W^\star_\varepsilon)$ and
$q_t=\E_\varepsilon[(P_t^\star)^2\mid\mathscr F_t]$.
Condition~\ref{ec:pricing} implies
\[
\E_\varepsilon[P_t^\star\mid\mathscr F_t]=q_t,\qquad 0\leq q_t\leq1,
\qquad
\E_\varepsilon[(m^\star_{t+1,\varepsilon})^4\mid\mathscr F_t]
\leq8(1+K_4).
\]

\item The score $m^\star_{t+1,\varepsilon}X_t$ is a square-integrable
Hilbert-valued martingale difference with respect to the economic timing:
\[
\E_\varepsilon[m^\star_{t+1,\varepsilon}X_t\mid\mathscr F_t]=0,
\qquad
\E_\varepsilon[(m^\star_{t+1,\varepsilon})^2 X_t\otimes X_t]
\preceq C_{\rm sc}\Sigma.
\]
\end{enumerate}
\end{lemma}

\begin{proof}
By the reproducing property,
\[
\|X_t\|_{\Hh_K}^2
=
\frac1{N^2}
\sum_{i,j=1}^N
R_{i,t+1}R_{j,t+1}K(Z_{i,t},Z_{j,t})
=
\frac1{N^2}R_{t+1}'K_tR_{t+1}.
\]
Hence Condition~\ref{ec:returns} gives part~(i) directly. Markov's
inequality yields the conditional tail bound, and the nonnegative power
series of the exponential yields the moment bound after taking expectations.
Measurability follows from the measurable kernel sections and separability.
Conditional Jensen followed by ordinary expectation gives
$\E\|X_t\|^2\leq B_X^2\log2$. Taking the expectation of the positive
rank-one operator in trace norm proves part~(ii).

Square integrability justifies differentiation of the population quadratic
criterion. Its first-order condition is
$\langle V,\mu-\Sigma W^\star_\varepsilon\rangle=0$ for every $V\in\Hh_K$.
Expanding around $W^\star_\varepsilon$ removes the cross term and proves
part~(iii). No bounded-support argument is used in these identities.

Taking $V=W^\star_\varepsilon$ in Condition~\ref{ec:pricing} yields
$\E[P_t^\star\mid\mathscr F_t]=q_t$.
Conditional Cauchy--Schwarz gives $q_t^2\leq q_t$, hence $0\leq q_t\leq1$.
The conditional fourth-moment bound in Condition~\ref{ec:returns} applies
to $P_t^\star$, since its stock weights are $\mathscr F_t$-measurable.
Thus
\[
\E[(1-P_t^\star)^4\mid\mathscr F_t]
\leq8\{1+\E[(P_t^\star)^4\mid\mathscr F_t]\}
\leq8(1+K_4q_t^2)\leq8(1+K_4).
\]
For every $V\in\Hh_K$, the same condition and conditional
Cauchy--Schwarz give
\[
\begin{aligned}
\E[(m^\star)^2(R^p(V))^2\mid\mathscr F_t]
&\leq\sqrt{\E[(m^\star)^4\mid\mathscr F_t]
                 \E[(R^p(V))^4\mid\mathscr F_t]}\\
&\leq C_{\rm sc}\E[(R^p(V))^2\mid\mathscr F_t].
\end{aligned}
\]
Integration proves the operator inequality in part~(v).
Summing over an orthonormal basis gives
$\E\|m^\star X_t\|^2\leq C_{\rm sc}\Tr\Sigma<\infty$.
The conditional expectation therefore exists. Taking its inner products
against a countable dense subset of $\Hh_K$ proves the martingale-difference
identity from Condition~\ref{ec:pricing}.
\end{proof}

\begin{proposition}[A sufficient Gaussian return specification]
\label{prop:gaussian_admissible}
Suppose that, conditionally on $\mathscr F_t$, $R_{t+1}$ is Gaussian, possibly
with a singular covariance, and
\[
\max_{1\leq i\leq N}\E[R_{i,t+1}^2\mid\mathscr F_t]\leq V_R<\infty
\quad\text{a.s.},
\]
uniformly in the economy and time. Then Condition~\ref{ec:returns} holds
with $B_X^2=8\kappa_KV_R$ and $K_4=3$.
\end{proposition}
\begin{proof}
For every $V\in\Hh_K$, $R^p_{t+1}(V)$ is conditionally Gaussian. A scalar
$G\sim N(m,v)$ satisfies
\[
\E G^4=m^4+6m^2v+3v^2\leq3(m^2+v)^2,
\]
which proves \eqref{eq:conditional_fourth_condition} with $K_4=3$.

For the tail condition, Cauchy--Schwarz in $\Hh_K$ and
$K(z,z)\leq\kappa_K$ give
\[
\|X_t\|_{\Hh_K}^2
\leq
\frac{\kappa_K}{N}\sum_{i=1}^N R_{i,t+1}^2.
\]
If $G\sim N(m,v)$ and $m^2+v\leq V_R$, then
\[
\E e^{G^2/(8V_R)}
=\left(1-\frac{v}{4V_R}\right)^{-1/2}
  \exp\!\left(\frac{m^2}{8V_R-2v}\right)
\leq\frac{2}{\sqrt3}e^{1/6}<2.
\]
Therefore, by convexity,
\[
\begin{aligned}
\E_\varepsilon\!\left[
\left.
\exp\!\left(\frac{\|X_t\|^2}{8\kappa_KV_R}\right)
\right|\mathscr F_t
\right]
&\leq
\E_\varepsilon\!\left[
\left.
\exp\!\left(\frac1N\sum_{i=1}^N
              \frac{R_{i,t+1}^2}{8V_R}\right)
\right|\mathscr F_t
\right]\\
&\leq
\frac1N\sum_{i=1}^N
\E_\varepsilon\!\left[
\left.
\exp\!\left(\frac{R_{i,t+1}^2}{8V_R}\right)
\right|\mathscr F_t
\right]
<2.
\end{aligned}
\]
Independence across stocks is unnecessary.
\end{proof}

The proposition is a sufficient condition, not a claim that every Gaussian
mixture has uniform covariance-relative sub-Gaussian tails. Conditional
Gaussianity with bounded conditional second moments is enough here.
No Bernstein condition is imposed after conditioning on the entire
augmented input $(Z_t,R_{t+1})$: at that point $m^\star$ is already known.

\subsection{Exact reduction to the kernel learning problem}
\label{app:induced_krr}

\begin{lemma}[Induced KRR problem]
\label{lem:induced_krr}
For every $\varepsilon\in\mathcal E$, the direct estimator is isometrically
equivalent to scalar KRR with constant response one. Under this identification,
the statistical covariance operator is $\Sigma_{\Hh_K,\varepsilon}$ and
the KRR prediction error is exactly population portfolio regret. The source
order is $r=1$. Since the induced kernel need not be uniformly bounded,
Theorem~3.2 of \citet{bignozzigagliardinischneider2026} is not invoked
verbatim; Lemma~\ref{lem:unbounded_stability} supplies the replacement
empirical-stability and score bounds.
\end{lemma}

\begin{proof}
\textit{Step 1: data and payoff representation.}
Let
\[
\mathcal X:=\mathscr Z^N\times\mathbb R^N,
\qquad
\Xi_t:=(Z_t,R_{t+1}).
\]
For $x=((z_1,\ldots,z_N),r)\in\mathcal X$, put
\[
\Phi(x):=\frac1N\sum_{i=1}^N r_iK(\cdot,z_i),
\qquad
f_W(x):=\frac1N\sum_{i=1}^N W(z_i)r_i.
\]
Then $\Phi(\Xi_t)=X_t$ and
$f_W(\Xi_t)=\langle W,X_t\rangle_{\Hh_K}=R^p_{t+1}(W)$.
The symbol $\Xi_t$ denotes the augmented observation; $X_t$ continues to
denote the managed-payoff feature used in the main text.

Let $\mathcal G:=\{f_W:W\in\Hh_K\}$ and equip it with
$\langle f_W,f_V\rangle_{\mathcal G}:=\langle W,V\rangle_{\Hh_K}$.
This is well defined: if $f_W$ vanishes on $\mathcal X$, choose
$r=(N,0,\ldots,0)$ and vary $z_1$ to obtain $W(z_1)=0$.
Thus the map $\mathcal T:W\mapsto f_W$ is a Hilbert-space isometry onto
$\mathcal G$. This uses the declared input domain, not a full-support
assumption on the actual return law.

With $Y_t\equiv1$, the empirical objectives agree exactly:
\[
\frac1T\sum_{t=1}^T|Y_t-f_W(\Xi_t)|^2+\lambda\|f_W\|_{\mathcal G}^2
=
\frac1T\sum_{t=1}^T(1-R^p_{t+1}(W))^2+\lambda\|W\|_{\Hh_K}^2.
\]
This is an algebraic reformulation of the portfolio criterion, not an
additional specification of the return-generating process.

\medskip
\noindent\textit{Step 2: induced kernel and population operator.}
The reproducing kernel of $\mathcal G$ is
\[
\widetilde K(x,x')
=\langle\Phi(x),\Phi(x')\rangle_{\Hh_K}
=\frac1{N^2}\sum_{i,j}r_i s_jK(z_i,z_j'),
\qquad x'=((z_1',\ldots,z_N'),s).
\]
It is measurable and $\widetilde K(x,x)=\|\Phi(x)\|^2$, without a uniform
bound on this diagonal. Nevertheless,
$\E_\varepsilon\widetilde K(\Xi_t,\Xi_t)<\infty$ by
Lemma~\ref{lem:primitive_consequences}, so the canonical embedding is bounded.
Let $J:\mathcal G\to L^2(\mathbb P_{\varepsilon,\Xi})$ be the canonical
embedding. For $W,V\in\Hh_K$,
\[
\langle J\mathcal TW,J\mathcal TV\rangle_{L^2}
=\E_\varepsilon[\langle W,X_t\rangle\langle V,X_t\rangle]
=\langle W,\Sigma_{\Hh_K,\varepsilon}V\rangle_{\Hh_K}.
\]
Consequently,
\[
\mathcal T^*J^*J\mathcal T=\Sigma_{\Hh_K,\varepsilon},
\qquad
\mathcal T^*J^*1=\mu_{\Hh_K,\varepsilon}.
\]
The minimum-norm population target is
$f_{\mathcal G}=\mathcal TW^\star_\varepsilon$, so
$\|f_{\mathcal G}\|_{\mathcal G}\leq R$. This is precisely source order
$r=1$. The positive eigenvalues and effective dimension are unchanged by
the isometry.

\medskip
\noindent\textit{Step 3: dependence, filtration, and statistical residual.}
The space $\mathcal X$ is Polish, and
$\{(\Xi_t,Y_t)\}$ inherits stationarity and the mixing bound from
Condition~\ref{ec:dependence}. Define its statistical history by
\[
\mathcal A_t:=\sigma(\Xi_s,Y_s:s\leq t),
\qquad
\mathcal B_t:=\mathcal A_{t-1}\vee\sigma(\Xi_t).
\]
The economic timing in Definition~\ref{def:economy} implies
$\mathcal A_{t-1}\subseteq\mathscr F_t$.
The statistical residual is
\[
\xi_t:=Y_t-f_{\mathcal G}(\Xi_t)
=1-R^p_{t+1}(W^\star_\varepsilon)
=m^\star_{t+1,\varepsilon}.
\]
By Lemma~\ref{lem:primitive_consequences}(v) and the tower property,
\[
\E_\varepsilon[X_t\xi_t\mid\mathcal A_{t-1}]=0.
\]
Since $\widetilde K(\cdot,\Xi_t)=\mathcal TX_t$, the corresponding
kernel-weighted residual in $\mathcal G$ is also a filtered innovation.
Lemma~\ref{lem:primitive_consequences} supplies the score variance bound.
The residual $\xi_t$ is a projection residual and is $\mathcal B_t$-measurable.
Consequently, an almost-sure conditional Bernstein bound given $\mathcal B_t$
would force it to be bounded; it is not imported from the companion theorem.
We also do not assume $\E[\xi_t\mid\Xi_t]=0$. The exact replacement for the
probabilistic step is given next.

\medskip
\noindent\textit{Step 4: identify the risk.}
For any fixed $W\in\Hh_K$, Lemma~\ref{lem:primitive_consequences}(iii)
gives
\[
\|Jf_W-Jf_{\mathcal G}\|_{L^2}^2
=\E_\varepsilon[(R^p_{t+1}(W)-R^p_{t+1}(W^\star_\varepsilon))^2]
=Q_\varepsilon(W)-Q_\varepsilon(W^\star_\varepsilon).
\]
The same identity applies to an estimated policy when the training data
are held fixed in the population evaluation. All constants are inherited
from the fixed representation and the parameters of $\mathcal E$.
\end{proof}

\subsection{Empirical stability with unbounded managed features}
\label{app:unbounded_concentration}

This subsection replaces the empirical-stability and noise proposition in
\citet{bignozzigagliardinischneider2026}. The covariance calculation uses
their alternating-block/Berbee argument, with an analytical truncation of
the feature norm. The score is controlled by a direct martingale second-moment
argument. The estimator itself is not truncated, and the remaining
operator decomposition and source-bias calculation are unchanged.

\begin{lemma}[Unbounded-feature stability and score]
\label{lem:unbounded_stability}
Assume Conditions~\ref{ec:dependence}--\ref{ec:pricing} and
$\Tr\Sigma<\infty$. Put
\[
C_{\rm sc}:=\sqrt{8K_4(1+K_4)}.
\]
For $0<\lambda\leq1$, put
\[
D_\lambda=(\Sigma+\lambda I)^{-1/2},\quad
C_\lambda=D_\lambda\Sigma D_\lambda,\quad
 d_\lambda=\Tr C_\lambda=\mathcal C_K(\lambda),
\]
\[
\Delta_{\lambda,T}=D_\lambda(\Sigma-\widehat\Sigma_T)D_\lambda,
\qquad
S_{\lambda,T}=\frac1T\sum_{t=1}^T D_\lambda m^\star_{t+1}X_t.
\]
Fix $\delta\in(0,1)$ and define
\[
\begin{split}
\ell&=\max\!\left\{1,\left\lceil
  (8C_\beta T/\delta)^{1/(\gamma+1)}\right\rceil\right\},\\
H&=\log\!\left(\frac{eT}{\delta\lambda}\right),\qquad
\tau^2=8B_X^2H,\qquad B_\lambda=\frac{1+\tau^2}{\lambda},\\
z&=\log\!\left(\frac{112(1+d_\lambda)}{\delta}\right).
\end{split}
\]
If $\ell\leq T/4$, there is a numerical constant $C$ such that, with
probability at least $1-\delta$,
\begin{equation}
\|\Delta_{\lambda,T}\|_{\op}
\leq C\!\left\{\sqrt{\frac{\ell B_\lambda z}{T}}
                  +\frac{\ell B_\lambda z}{T}\right\}
       +\frac{2B_X^2(8H+1)}{\lambda}e^{-8H},
\label{eq:unbounded_covariance}
\end{equation}
and
\begin{equation}
\|S_{\lambda,T}\|_{\Hh_K}^2
\leq\frac{2C_{\rm sc}}{\delta}\frac{d_\lambda}{T}.
\label{eq:unbounded_score}
\end{equation}
Consequently, for each fixed $\delta$ there is a class-uniform $C_\delta$
such that, for sufficiently large $T$, the sufficient condition
\begin{equation}
T^{\gamma/(\gamma+1)}\lambda
\geq C_\delta\left[\log\!\left(\frac{C_\delta T}{\lambda}\right)\right]^2
\label{eq:unbounded_admissibility}
\end{equation}
implies $\|\Delta_{\lambda,T}\|_{\op}\leq1/2$ and
\eqref{eq:unbounded_score} with probability at least $1-\delta$.
\end{lemma}

\begin{proof}
We suppress the economy index. Only the covariance is truncated in the
proof. Let
\[
A_t=(D_\lambda X_t)\otimes(D_\lambda X_t)
       \mathbf1_{\{\|X_t\|\leq\tau\}},\qquad
C_{\lambda,\tau}=\E A_t,\qquad H_t=A_t-C_{\lambda,\tau}.
\]
The process $(H_t)$ is centered and has no larger absolute-regularity
coefficients than $(Z_t,R_{t+1})$. Since
$0\preceq C_{\lambda,\tau}\preceq C_\lambda\preceq I$ and
$\|A_t\|_{\op}\leq\tau^2/\lambda$, we have
\begin{equation}
\|H_t\|_{\op}\leq B_\lambda,\qquad
\E H_t^2=\E A_t^2-C_{\lambda,\tau}^2
\preceq B_\lambda C_\lambda.
\label{eq:truncated_operator_variance}
\end{equation}
The second inequality is an operator inequality; no independence of
successive observations is used.

Let $n=\lfloor T/(2\ell)\rfloor$ and split the first $2n\ell$
observations into $2n$ consecutive blocks of length $\ell$.
Treat the odd and even blocks separately. For an odd-block sum $B_j^o$,
operator Cauchy--Schwarz gives
\[
\|B_j^o\|_{\op}\leq\ell B_\lambda,\qquad
\E(B_j^o)^2\preceq\ell^2B_\lambda C_\lambda.
\]
Indeed, for every $h$, $\|\sum_{t\in I}H_th\|^2
\leq\ell\sum_{t\in I}\|H_th\|^2$.
Berbee coupling replaces each family by independent blocks with identical
block laws. The total probability of a coupling failure for both families is
at most
\[
2n\beta(\ell)\leq C_\beta T\ell^{-(\gamma+1)}\leq\delta/8.
\]
For the independent copies, a variance majorant has operator norm at most
$n\ell^2B_\lambda$ and trace at most $n\ell^2B_\lambda d_\lambda$.
The trace form of the self-adjoint Bernstein inequality
\citep[proof of Theorem~3.1 and Section~3.2]{minsker2017}, with this
variance majorant, yields for each family, with failure probability at most
$\delta/8$,
\[
\left\|\frac1T\sum_{j=1}^n B_j^{o,*}\right\|_{\op}
\leq C\!\left\{\sqrt{\frac{\ell B_\lambda z}{T}}
                       +\frac{\ell B_\lambda z}{T}\right\}.
\]
The trace-to-variance-majorant ratio is at most $d_\lambda$; no bound on
the effective rank of the actual variance is inferred by monotonicity.
Apply the finite-dimensional inequality to increasing finite-rank
projections and pass to the limit in operator norm, as in that reference.
The operators here are trace class and bounded after truncation.
The fewer than $2\ell$ leftover observations contribute at most
$2\ell B_\lambda/T$, already absorbed because $z\geq1$.

It remains to compare the truncated and original covariances.
Lemma~\ref{lem:primitive_consequences} gives
\[
\mathbb P\!\left(\max_{t\leq T}\|X_t\|>\tau\right)
\leq2Te^{-8H}\leq\delta/8.
\]
On the complementary event, the truncated empirical covariance equals the
original empirical covariance. The population truncation bias satisfies
\[
\begin{aligned}
\|C_\lambda-C_{\lambda,\tau}\|_{\op}
&\leq\lambda^{-1}\E[\|X_t\|^2\mathbf1_{\{\|X_t\|>\tau\}}]\\
&\leq\frac{2}{\lambda}(\tau^2+B_X^2)e^{-\tau^2/B_X^2}
=\frac{2B_X^2(8H+1)}{\lambda}e^{-8H}.
\end{aligned}
\]
The tail-integration formula justifies the second line. This proves
\eqref{eq:unbounded_covariance} with total failure probability at most
$\delta/2$.

For the score, Lemma~\ref{lem:primitive_consequences} gives
\[
\E\|D_\lambda m^\star_{t+1}X_t\|^2
\leq C_{\rm sc}\Tr(D_\lambda\Sigma D_\lambda)
=C_{\rm sc}d_\lambda.
\]
For $s<t$, the score at $s$ is $\mathscr F_t$-measurable, whereas the
conditional mean of the score at $t$ given $\mathscr F_t$ is zero.
All cross inner products consequently have zero expectation. Hence
\[
\E\|S_{\lambda,T}\|^2
=\frac1{T^2}\sum_{t=1}^T
  \E\|D_\lambda m^\star_{t+1}X_t\|^2
\leq C_{\rm sc}\frac{d_\lambda}{T}.
\]
Markov's inequality proves \eqref{eq:unbounded_score} with failure
probability $\delta/2$. There is no blocking loss in this score bound.

Finally, $d_\lambda\leq B_X^2\log2/\lambda$ and, for fixed $\delta$,
$\ell\leq C_\delta T^{1/(\gamma+1)}$.
Thus the leading square-root term in \eqref{eq:unbounded_covariance} is
bounded by a constant times
\[
\left\{
\frac{[\log(C_\delta T/\lambda)]^2}
     {T^{\gamma/(\gamma+1)}\lambda}
\right\}^{1/2}.
\]
The linear term is bounded by the corresponding squared expression.
The truncation bias tends to zero uniformly for $0<\lambda\leq1$ as
$T\to\infty$. Increasing $C_\delta$ proves the final statement.
\end{proof}

\begin{corollary}[Finite-sample portfolio loss]
\label{cor:unbounded_risk}
Under the assumptions of Lemma~\ref{lem:unbounded_stability}, for each fixed
$\delta\in(0,1)$ and sufficiently large $T$ satisfying
\eqref{eq:unbounded_admissibility},
\[
Q_\varepsilon(\widehat W_{\lambda,T})-Q_\varepsilon(W^\star_\varepsilon)
\leq9R^2\lambda+C_\delta\frac{\mathcal C_K(\lambda)}{T}
\]
with probability at least $1-\delta$, uniformly over the class.
\end{corollary}
\begin{proof}
Use the error decomposition, source-bias bound and operator-interaction
bound of \citet{bignozzigagliardinischneider2026}, specialized to $r=1$.
For completeness, write
\[
\begin{split}
W_\lambda&=(\Sigma+\lambda I)^{-1}\Sigma W^\star,\qquad B=W_\lambda-W^\star,\\
V_1&=(\widehat\Sigma_T+\lambda I)^{-1}
          \frac1T\sum_t m^\star_{t+1}X_t,\\
V_2&=(\widehat\Sigma_T+\lambda I)^{-1}(\Sigma-\widehat\Sigma_T)B.
\end{split}
\]
Then $\widehat W-W^\star=B+V_1+V_2$.
Spectral calculus gives $\|\Sigma^{1/2}B\|^2\leq R^2\lambda$ and
$\|(\Sigma+\lambda I)^{1/2}B\|^2\leq2R^2\lambda$.
On the stability event,
\[
\|\Sigma^{1/2}V_1\|^2\leq4\|S_{\lambda,T}\|^2,
\qquad
\|\Sigma^{1/2}V_2\|^2\leq2R^2\lambda.
\]
The first inequality follows from $\|(I-\Delta_{\lambda,T})^{-1}\|\leq2$;
the second follows from
$\|(I-\Delta_{\lambda,T})^{-1}\Delta_{\lambda,T}\|\leq1$.
Finally, $\|u+v+w\|^2\leq3(\|u\|^2+\|v\|^2+\|w\|^2)$
and Lemma~\ref{lem:unbounded_stability} give the claim. In particular,
one may take the score constant $24C_{\rm sc}/\delta$.
\end{proof}


\subsection{From portfolio loss to the Sharpe ratio}
\label{app:sharpe_recovery}

\begin{lemma}[Sharpe-ratio recovery without loss of exponent]
\label{lem:sharpe_recovery}
Let $\widehat W_T\in\Hh_K$ be any estimated policy satisfying, uniformly over $\varepsilon\in\mathcal E$,
\[
D_T:=Q_\varepsilon(\widehat W_T)-Q_\varepsilon(W^\star_\varepsilon)
=O_{\mathbb P_\varepsilon}(a_T),
\qquad a_T\longrightarrow0.
\]
Then the estimated payoff has positive mean and positive variance with
probability tending to one, uniformly, and
\[
\SR^\star_{\varepsilon,K}
-\SR(R^p_{t+1}(\widehat W_T))
=O_{\mathbb P_\varepsilon}(a_T).
\]
Sharpe may be assigned any finite convention on a zero-variance event,
whose probability tends to zero. No sign adjustment using unknown
population moments is required.
\end{lemma}

\begin{proof}
Write $q^\star=Q_\varepsilon(W^\star_\varepsilon)$ and
$S^\star=\SR^\star_{\varepsilon,K}$. By
Proposition~\ref{prop:sdf_hj} and Condition~\ref{ec:sharpe},
\[
\underline q:=\frac1{1+\overline S^2}
\leq q^\star\leq
\frac1{1+\underline S^2}=: \overline q<1.
\]

\medskip
\noindent\textit{Step 1: recover the sign and exclude degeneracy.}
On the event $D_T<1-\overline q$, which has probability tending to one
uniformly, $Q_\varepsilon(\widehat W_T)<1$. Expanding this loss yields
\[
2\E_\varepsilon[R^p_{t+1}(\widehat W_T)]
>
\E_\varepsilon[(R^p_{t+1}(\widehat W_T))^2]\geq0.
\]
Thus its mean is strictly positive. It cannot have zero variance: a zero-variance payoff with positive mean is a positive constant almost surely, and scaling the same policy would produce the unit payoff, contradicting $q^\star\geq\underline q>0$. Hence its Sharpe ratio
$\widehat S_T$ is positive on this event.

\medskip
\noindent\textit{Step 2: bound squared-Sharpe loss.}
Optimizing only the scale of $\widehat W_T$ gives
\[
q^\star
\leq\frac1{1+\widehat S_T^2}
=\min_{a\in\mathbb R}Q_\varepsilon(a\widehat W_T)
\leq Q_\varepsilon(\widehat W_T)
=q^\star+D_T.
\]
Therefore
\[
0\leq(S^\star)^2-\widehat S_T^2
\leq\frac1{q^\star}-\frac1{q^\star+D_T}
=\frac{D_T}{q^\star(q^\star+D_T)}
\leq\frac{D_T}{\underline q^2}.
\]

\medskip
\noindent\textit{Step 3: pass to the level Sharpe.}
Since $S^\star+\widehat S_T\geq\underline S$ on the same event,
\[
0\leq S^\star-\widehat S_T
=\frac{(S^\star)^2-\widehat S_T^2}{S^\star+\widehat S_T}
\leq\frac{D_T}{\underline S\,\underline q^2}.
\]
This proves the stated rate uniformly over $\mathcal E$. The complementary event has uniformly vanishing probability.
\end{proof}

\subsection{Proof of the portfolio-learning rate}
\label{app:rate}

\begin{proof}[Proof of Theorem~\ref{thm:rate}]
Fix $\varepsilon\in\mathcal E$.

\medskip
\noindent\textit{Step 1: use the unbounded-feature extension.}
The proof follows the operator decomposition of
\citet{bignozzigagliardinischneider2026}, with its empirical-stability and
noise step replaced by Lemma~\ref{lem:unbounded_stability}.
Corollary~\ref{cor:unbounded_risk} therefore gives, for every fixed
$\delta\in(0,1)$, a constant $C_\delta$ common to the class such that
\[
Q_\varepsilon(\widehat W_{\lambda,T})-Q_\varepsilon(W^\star_\varepsilon)
\leq C_\delta\left(\lambda+\frac{\mathcal C_{K,\varepsilon}(\lambda)}{T}\right)
\]
with probability at least $1-\delta$, for sufficiently large $T$ satisfying
\begin{equation}
T^{\gamma/(\gamma+1)}\lambda
\geq C_\delta[\log(C_\delta T/\lambda)]^2.
\label{eq:sample_stability}
\end{equation}
The extra logarithm occurs only in this sufficient stability condition,
not in the leading score variance. By Condition~\ref{ec:spectrum} and
Proposition~A.6 of \citet{bignozzigagliardinischneider2026},
$\mathcal C_{K,\varepsilon}(\lambda)\asymp\lambda^{-1/b}$ uniformly.

\medskip
\noindent\textit{Step 2: verify the admissible regularization scale.}
For $\lambda_T^\star\asymp T^{-b/(b+1)}$, the class restriction
$\gamma>b$ implies
\[
\frac{\gamma}{\gamma+1}-\frac{b}{b+1}
=\frac{\gamma-b}{(\gamma+1)(b+1)}>0.
\]
Consequently,
\[
\frac{T^{\gamma/(\gamma+1)}\lambda_T^\star}
 {[\log(C_\delta T/\lambda_T^\star)]^2}\longrightarrow\infty,
\]
so \eqref{eq:sample_stability} holds eventually. At this scale,
\[
\lambda_T^\star
\asymp\frac{\mathcal C_{K,\varepsilon}(\lambda_T^\star)}{T}
\asymp T^{-b/(b+1)}.
\]
This proves the population-loss rate. Uniformity means that, for each
$\delta\in(0,1)$, some $C_\delta,T_0(\delta)$ independent of $\varepsilon$
satisfy
\[
\sup_{\varepsilon\in\mathcal E}
\mathbb P_\varepsilon\!\left(
 Q_\varepsilon(\widehat W_{\lambda_T^\star,T})
 -Q_\varepsilon(W^\star_\varepsilon)
 >C_\delta T^{-b/(b+1)}
\right)\leq\delta,
\qquad T\geq T_0(\delta).
\]

\medskip
\noindent\textit{Step 3: translate loss into Sharpe and frontier recovery.}
Apply Lemma~\ref{lem:sharpe_recovery} with $a_T=T^{-b/(b+1)}$. At any fixed volatility target $\sigma$, the difference
between the population and learned capital-allocation lines is
\[
\mu^\star_K(\sigma)-\widehat\mu_{K,T}(\sigma)
=\sigma\left(\SR^\star_{\varepsilon,K}
-\SR(R^p_{t+1}(\widehat W_{\lambda_T^\star,T}))\right),
\]
which has the same order, also uniformly for $\sigma$ in a bounded interval.
This concerns the tangency capital-allocation line with a risk-free asset
and scalar rescaling, not the entire risky-only frontier.
\end{proof}

\subsection{Effective portfolio complexity}
\label{app:complexity}

\begin{proof}[Proof of Theorem~\ref{thm:effective_complexity}]
The proof of Theorem~\ref{thm:rate} already establishes admissibility of
$\lambda_T^\star$ and the spectral relation
$\mathcal C_{K,\varepsilon}(\lambda)\asymp\lambda^{-1/b}$. Thus
\[
\mathcal C^\star_{K,T}
=\mathcal C_{K,\varepsilon}(\lambda_T^\star)
\asymp T^{1/(b+1)},
\qquad
\frac{\mathcal C^\star_{K,T}}{T}
\asymp T^{-b/(b+1)}\longrightarrow0.
\]
It is finite for every $T$, since
$\mathcal C_{K,\varepsilon}(\lambda)\leq
\Tr(\Sigma_{\Hh_K,\varepsilon})/\lambda$.

To express the same balance in complexity units, let
$C=\mathcal C_{K,\varepsilon}(\lambda)$, so that $\lambda\asymp C^{-b}$.
The leading risk upper envelope has the form
\[
\mathfrak R_T(C)=aC^{-b}+d\frac CT,
\qquad a,d>0.
\]
Its derivatives are
\[
\mathfrak R_T'(C)=-baC^{-b-1}+\frac dT,
\qquad
\mathfrak R_T''(C)=b(b+1)aC^{-b-2}>0.
\]
Hence its unique minimum occurs at
\[
C_T^{\mathrm{env}}
=\left(\frac{ba}{d}T\right)^{1/(b+1)}.
\]
This is an upper-envelope calculation, not a shape theorem for the exact
Sharpe curve of every economy. No matching pointwise bias lower bound is used. Condition~\ref{ec:sharpe} is not needed for this conclusion.
\end{proof}


\subsection{A finance-specific minimax construction}
\label{app:minimax}

\begin{proof}[Proof of Theorem~\ref{thm:minimax_words}]
The construction keeps the policy kernel fixed. It is a least-favourable family
used only in this lower-bound proof, not a structural return equation imposed on
the economies in the main text.

\medskip
\noindent\textit{Step 1: construct a least-favourable family with nondegenerate Sharpe.}
Since $\mathcal E$ is nonempty, take the one-period state-return marginal
$\nu_0$ of one economy in the class. Write a draw as
$x=(\mathbf z,r)$ and put
\[
Q_X(x):=\|\Phi(x)\|_{\Hh_K}^2,
\qquad
q_X:=\int Q_X(x)\,d\nu_0(x).
\]
Condition~\ref{ec:returns} and Jensen's inequality imply
$q_X\leq B_X^2\log2$. Moreover $q_X>0$, because otherwise the managed
feature would vanish almost surely, contradicting the spectral lower bound
in Condition~\ref{ec:spectrum}.

Define a tilted law by
\[
d\widetilde\nu(x)
=
\frac{Q_X(x)}{q_X}\,d\nu_0(x),
\]
and, on the support of the tilted law, replace the return vector by
\[
r^\circ
=
\sqrt{\frac{q_X}{Q_X(x)}}\,r.
\]
Let $\nu$ denote the law of $(\mathbf z,r^\circ)$ and reuse $r$ for the
normalized return vector below. Since the feature map is linear in returns,
\[
\int \Phi(\mathbf z,r^\circ)\otimes\Phi(\mathbf z,r^\circ)\,d\widetilde\nu
=
\int \Phi(x)\otimes\Phi(x)\,d\nu_0(x)
=:\Sigma_0.
\]
Thus the normalization preserves the entire managed-payoff second-moment
operator. Under $\nu$,
\[
\|\Phi(x)\|_{\Hh_K}^2=q_X\leq B_X^2\log2
\quad\text{almost surely}.
\]
Put $\overline B_X:=B_X\sqrt{\log2}$.

Write
\[
\Sigma_0e_j=\mu_je_j,
\qquad
c_\mu j^{-b}\leq\mu_j\leq C_\mu j^{-b}.
\]
Its eigenvectors are orthonormal in $\Hh_K$. Fix constants $a,h>0$,
with $h\leq a$, sufficiently small that
\[
a^2+h^2\leq R^2,
\qquad
\overline B_X\sqrt{a^2+h^2}\leq\frac12,
\qquad
8C_\mu h^2\leq\frac18.
\]
For an integer $M\geq1$ and $\theta\in\{-1,1\}^M$, define
\[
W_\theta
:=
a e_1
+
\frac{h}{\sqrt M}
\sum_{j=M+1}^{2M}\theta_{j-M}e_j,
\qquad
g_\theta(x)
:=
\langle W_\theta,\Phi(x)\rangle_{\Hh_K}.
\]
Then
\[
\|W_\theta\|_{\Hh_K}=\sqrt{a^2+h^2}\leq R,
\qquad
|g_\theta(x)|\leq \overline B_X\sqrt{a^2+h^2}\leq\frac12.
\]
Put
\[
A_M
:=
\|g_\theta\|_{L^2(\nu)}^2
=
a^2\mu_1
+
\frac{h^2}{M}
\sum_{j=M+1}^{2M}\mu_j.
\]
This quantity does not depend on $\theta$. Moreover,
\[
a^2c_\mu\leq A_M\leq\frac14.
\]

\medskip
\noindent\textit{Step 2: construct admissible financial observations.}
Draw $x_t=(Z_t,r_t)$ independently with law $\nu$. Conditional on $x_t$,
draw $s_t\in\{-1,1\}$ according to
\[
\mathbb P_\theta(s_t=1\mid x_t)
=
\frac{1+g_\theta(x_t)}{2}.
\]
The observed financial row is
$(Z_t,R_{t+1})=(Z_t,s_tr_t)$. For the least-favourable construction take
\[
\mathscr F_t
=
\sigma(x_u,s_u:u<t)\vee\sigma(x_t).
\]
This information structure is allowed by Definition~\ref{def:economy}.
Rows are i.i.d. The managed feature is $s_t\Phi(x_t)$ and, since
$s_t^2=1$,
\[
\|s_t\Phi(x_t)\|_{\Hh_K}^2=q_X.
\]
Because $x_t$ is contained in $\mathscr F_t$,
\[
\E_\theta\!\left[
\left.
\exp\!\left(
\frac{\|s_t\Phi(x_t)\|_{\Hh_K}^2}{B_X^2}
\right)
\right|\mathscr F_t
\right]
=
\exp(q_X/B_X^2)
\leq2.
\]
Moreover, for every $V\in\Hh_K$,
\[
(R^p_{t+1}(V))^2=\langle V,\Phi(x_t)\rangle_{\Hh_K}^2,
\qquad
(R^p_{t+1}(V))^4=\langle V,\Phi(x_t)\rangle_{\Hh_K}^4
\]
conditionally on $\mathscr F_t$. Hence
Condition~\ref{ec:returns} holds with fourth-moment coefficient one,
and therefore with the class constant $K_4\geq1$.
The second-moment operator of the managed feature is exactly $\Sigma_0$.

For every $V\in\Hh_K$, conditional on the current $x_t$ and the past,
\[
\begin{aligned}
&\E_\theta\!\left[
(1-R^p_{t+1}(W_\theta))R^p_{t+1}(V)
\mid x_t,\mathscr F_{t-1}
\right]\\
&\qquad=
\langle V,\Phi(x_t)\rangle_{\Hh_K}
\E_\theta[s_t-g_\theta(x_t)\mid x_t,\mathscr F_{t-1}]
=0.
\end{aligned}
\]
Thus Condition~\ref{ec:pricing} holds. For any portfolio policy $W$ with
finite loss, write $g_W(x)=N^{-1}\sum_iW(z_i)r_i$. Then
\begin{equation}
Q_\theta(W)-Q_\theta(W_\theta)
=
\int(g_W(x)-g_\theta(x))^2\,d\nu(x).
\label{eq:lower_risk_identity}
\end{equation}
Indeed, conditioning on $x$ gives the quadratic
$1-2g_\theta(x)g_W(x)+g_W(x)^2$. Consequently, $W_\theta$ is the
population-optimal minimum-norm policy.

Its payoff satisfies
\[
\E_\theta[R^p_{t+1}(W_\theta)]=A_M,
\qquad
\E_\theta[(R^p_{t+1}(W_\theta))^2]=A_M.
\]
Hence
\[
(\SR_\theta^\star)^2
=
\frac{A_M}{1-A_M}.
\]
Because $A_M$ is uniformly bounded away from zero and from one, the optimal
Sharpe ratio is uniformly bounded away from zero and infinity. Thus the
least-favourable family satisfies Condition~\ref{ec:sharpe}, with constants
common to the family, as well as Conditions~\ref{ec:dependence}--\ref{ec:spectrum}.

\medskip
\noindent\textit{Step 3: neighboring economies are difficult to distinguish.}
Consider the more informative experiment that also reveals $x_t$ and $s_t$.
A lower bound in this augmented experiment also applies to the financial
observations. Neighboring parameters differ in one sign, say coordinate
$j-M$. Their conditional Bernoulli probabilities lie in $[1/4,3/4]$, and
\[
\operatorname{KL}(\operatorname{Ber}(p),\operatorname{Ber}(q))
\leq
\frac{(p-q)^2}{q(1-q)}.
\]
Hence, for the laws $\mathbb P^{\mathrm{aug},T}_\theta$ of $T$ augmented
observations and a one-bit neighbor $\theta'$, we have
\[
\begin{aligned}
\operatorname{KL}(\mathbb P^{\mathrm{aug},T}_\theta,
                  \mathbb P^{\mathrm{aug},T}_{\theta'})
&\leq
2T\int(g_\theta-g_{\theta'})^2\,d\nu\\
&=
\frac{8Th^2\mu_j}{M}\\
&\leq
8C_\mu h^2\frac{T}{M^{b+1}}.
\end{aligned}
\]
Choose
\[
M=\left\lceil T^{1/(b+1)}\right\rceil.
\]
The right-hand side is at most $1/8$, so Pinsker's inequality gives total
variation at most $1/4$ for every neighboring pair.

\medskip
\noindent\textit{Step 4: testing difficulty implies scale-free directional error.}
Define
\[
\psi_j(x)
:=
\frac{\langle e_j,\Phi(x)\rangle_{\Hh_K}}{\sqrt{\mu_j}},
\]
so that $(\psi_j)_{j\geq1}$ is orthonormal in $L^2(\nu)$. Let $\widehat g$
be any estimated square-integrable payoff rule and write
\[
\widehat a_j
:=
\langle\widehat g,\psi_j\rangle_{L^2(\nu)}.
\]
Decode the signs by
\[
\widehat\theta_{j-M}
:=
\operatorname{sign}(\widehat a_1\widehat a_j),
\qquad M<j\leq2M,
\]
with any fixed convention at zero. This decoder is invariant to any nonzero
global rescaling of $\widehat g$.

Define the scale-free directional error
\[
D_\theta(\widehat g)
:=
\inf_{c\in\mathbb R}
\|g_\theta-c\widehat g\|_{L^2(\nu)}^2.
\]
For any $c\in\mathbb R$, if $c\widehat a_1\leq0$, the squared error in the
first coordinate is at least $a^2\mu_1$. If $c\widehat a_1>0$, every
incorrectly decoded sign contributes at least $h^2\mu_j/M$ to the squared
error. Since $h\leq a$ and $H(\widehat\theta,\theta)\leq M$, both cases
imply
\[
D_\theta(\widehat g)
\geq
\frac{h^2c_\mu}{2^bM^{b+1}}
H(\widehat\theta,\theta).
\]
For every neighboring pair, the sum of the two sign-error probabilities is at
least $1-\mathrm{TV}\geq3/4$. Averaging over the hypercube and summing over
its $M$ coordinates therefore gives
\[
2^{-M}\sum_\theta
\E_\theta[H(\widehat\theta,\theta)]
\geq
\frac{3M}{8}.
\]
Consequently,
\[
2^{-M}\sum_\theta
\E_\theta[D_\theta(\widehat g)]
\geq
\frac{3h^2c_\mu}{2^{b+3}}M^{-b}.
\]

\medskip
\noindent\textit{Step 5: directional error is a Sharpe-ratio gap.}
Write
\[
m:=\langle g_\theta,\widehat g\rangle_{L^2(\nu)},
\qquad
v:=\|\widehat g\|_{L^2(\nu)}^2.
\]
If $v>0$, the payoff generated by $\widehat g$ satisfies
\[
\E_\theta[R(\widehat g)]=m,
\qquad
\E_\theta[R(\widehat g)^2]=v,
\]
and hence
\[
\SR_\theta(\widehat g)^2
=
\frac{m^2}{v-m^2}.
\]
Moreover,
\[
D_\theta(\widehat g)
=
A_M-\frac{m^2}{v}
=
\frac{1}{1+\SR_\theta(\widehat g)^2}
-
\frac{1}{1+(\SR_\theta^\star)^2}.
\]
The same identity holds for $v=0$ under the convention
$\SR_\theta(\widehat g)=0$. Cauchy--Schwarz implies
$|\SR_\theta(\widehat g)|\leq\SR_\theta^\star$.

If $\SR_\theta(\widehat g)\geq0$, then
\[
\SR_\theta^\star-\SR_\theta(\widehat g)
\geq
\frac{D_\theta(\widehat g)}{2\overline S}.
\]
If $\SR_\theta(\widehat g)<0$, then the gap is at least $\underline S$;
since $D_\theta(\widehat g)\leq A_M\leq1/4$, it is again bounded below by
a fixed positive multiple of $D_\theta(\widehat g)$. Therefore there is a
constant $c_S>0$, independent of $M,T,\theta$, such that
\[
\SR_\theta^\star-\SR_\theta(\widehat g)
\geq
c_S D_\theta(\widehat g).
\]
Combining this inequality with Step 4 yields
\[
2^{-M}\sum_\theta
\E_\theta\!\left[
\SR_\theta^\star-\SR_\theta(\widehat g)
\right]
\geq
c_1M^{-b}
\]
for a constant $c_1>0$. Since
$M=\lceil T^{1/(b+1)}\rceil\leq2T^{1/(b+1)}$ for all sufficiently large
$T$, we obtain
\[
\sup_\theta
\E_\theta\!\left[
\SR_\theta^\star-\SR_\theta(\widehat g)
\right]
\geq
c_2T^{-b/(b+1)}
\]
for some $c_2>0$.

The argument allowed $\widehat g$ to range over all square-integrable
estimators in the more informative augmented experiment, a larger class than
managed portfolio rules based on the observed financial rows. Taking the
infimum over all measurable portfolio-learning procedures proves
\[
\inf_{\widetilde W_T}
\sup_{\varepsilon\in\mathcal E}
\mathbb E_\varepsilon
\left[
\SR^\star_{\varepsilon,K}
-
\SR_\varepsilon\!\left(R^p_{t+1}(\widetilde W_T)\right)
\right]
\geq
c\,T^{-b/(b+1)}
\]
for a constant $c>0$ and all sufficiently large $T$.
\end{proof}
